\subsection{Integral low-degree comparison for band-preserving quotients}

\begin{lemma}[Proved: integral projective-column bridge comparison]
\label{CP2:A-integral-column-bridge-connectivity}
Let $q:M\twoheadrightarrow H$ be a surjective homomorphism of
finite monoids.  Suppose $e\in M$ is idempotent, $f=q(e)$, and
the map of left $M$-sets $Me\to Hf$ induced by $q$ is bijective.
For any finite left-$N$, right-$M$ biset $S$, put
$U=S\otimes_M H$.  The actual map of two-sided bridges induced
by $q$ and $S\to U$ is an isomorphism on integral homology in
degrees zero and one and a surjection in degree two.
The same assertions hold with coefficients in any field.
\end{lemma}
\begin{proof}
Fix $k=\mathbb Z$ or a field.  Write $R=kM$, $B=kH$,
$T=kN$, $L=kS$, and $L'=kU$.  Set
\[
 F=L\otimes_R^{\mathbf L}B,\qquad
 D=\operatorname{fib}(F\longrightarrow L').
\]
The actual set coend gives $L'=L\otimes_R B$; hence
$H_i(D)=0$ for $i\leq0$.  All these complexes carry their
natural commuting left $T$ action.

The hypothesis identifies the left $R$ module $Bf$ with
the projective module $Re$.  Consequently
\[
 Ff=L\otimes_R^{\mathbf L}Bf\simeq Le
       =L\otimes_R Bf=L'f ,
\]
and the displayed comparison is the restriction of the actual
augmentation $F\to L'$.  It follows that $Df$ is acyclic.
Let $V_f=\ker(Bf\to k)$, where the map is augmentation.
The exact sequence
\[
       0\longrightarrow V_f\longrightarrow Bf
                      \longrightarrow k\longrightarrow0
\]
therefore gives
\[
 D\otimes_B^{\mathbf L}k
       \simeq (D\otimes_B^{\mathbf L}V_f)[1].
                                                               \tag{1}
\]
Here homological shift $[1]$ raises degrees by one.
Derived tensor products of complexes concentrated homologically
in degrees at least one and zero are concentrated in degrees
at least one.  Thus the complex in (1) has zero homology below
degree two.  Applying $k\otimes_T^{\mathbf L}-$ preserves this
connectivity.  Derived associativity identifies the resulting
comparison with the actual double-bar map
\[
 k\otimes_T^{\mathbf L}L\otimes_R^{\mathbf L}k
   \longrightarrow
 k\otimes_T^{\mathbf L}L'\otimes_B^{\mathbf L}k .
\]
Its fiber has zero homology below degree two, which gives
the asserted isomorphisms and surjection.

All tensor comparisons can be represented by the ordinary
nonnegative bar resolutions.  Their terms are free over the
appropriate monoid ring, since the monoid and biset modules
are free over $k$.  This supplies the needed connective models
over $\mathbb Z$ as well as over a field; no division or group
averaging occurs.
\end{proof}

\begin{theorem}[Proved: integral first and second homology of a band quotient]
\label{CP2:A-band-quotient-integral-low-degree}
Suppose $M$ has a proper rectangular-band minimal ideal $K$,
and $q:M\twoheadrightarrow H$ is bijective on $K$.
Identify $q(K)$ with $K$ using this bijection.
For every finite left-$N$, right-$M$ biset $S$, the actual
functor
\[
 \mathcal C(N,S,M)\longrightarrow
 \mathcal C(N,S\otimes_M H,H)
                                                               \tag{2}
\]
induces isomorphisms on integral homology in degrees zero,
one and two.  It induces an isomorphism on second cohomology
with any abelian coefficient group, in particular with
$\mathbb F_2$ coefficients.  The induced state set has at most
$|S|$ elements.
\end{theorem}
\begin{proof}
The image of the minimal ideal is the minimal ideal of $H$.
Choose $e\in K$.  Both $Me$ and $He$ are the corresponding
minimal columns, so the bijection on $K$ gives the
projective-column hypothesis of the preceding lemma.

We check the corner comparison integrally.  Set $R=\mathbb ZM$
and let
\[
 V=\ker(Re\longrightarrow\mathbb Z),\qquad
 W=\ker(eR\longrightarrow\mathbb Z).
\]
The two modules are the free row- and column-augmentation
modules.  Both are annihilated by $\mathbb ZK$:
a band element sends every row basis vector to the same
row basis vector, and the right-handed statement is dual.
The two projective augmentation sequences give naturally
\[
 H_2(BM;\mathbb Z)
   =\operatorname{Tor}^{R}_2(\mathbb Z,\mathbb Z)
   \cong\operatorname{Tor}^{R}_1(\mathbb Z,V)
   \cong W\otimes_R V.                                         \tag{3}
\]
For the second isomorphism, the tensor target
$eR\otimes_R V=eV$ is zero, so the low-degree end of
the right augmentation sequence gives exactly the last term
in (3).  The same argument applies to $H$.

The chosen band bijection identifies both augmentation
modules and their bases.  All their actions factor through
$q$, because multiplication by a band element is preserved.
Since $q$ is surjective, the balanced relations defining
$W\otimes_{\mathbb ZM}V$ and
$W\otimes_{\mathbb ZH}V$ coincide.  The natural map (3)
is therefore an isomorphism.  In addition, both corner
universal groups are trivial: every element of the band
maps to the identity in any group, and $m e\in K$ implies
that the image of every $m$ is the identity.  Hence the
first integral corner homology groups vanish.

Use the natural two-block homotopy-pushout models for (2).
Their corner maps are isomorphisms on $H_1$ and $H_2$,
and the opposite corner is unchanged.  The bridge map is
an isomorphism on $H_0,H_1$ and surjective on $H_2$ by
the preceding lemma.  Surjectivity identifies the images
of the two bridge-to-corner maps in degree two.  Thus their
cokernels are isomorphic; the adjacent kernels in degree
one are also isomorphic.  The two short exact sequences
extracted from Mayer--Vietoris give the asserted isomorphism
on $H_2$ of the collages.  The next portion of the same
sequence gives the $H_1$ isomorphism, and the degree-zero
comparison gives the $H_0$ isomorphism.

The natural integral universal-coefficient exact sequence
then gives the second-cohomology isomorphism with any
abelian coefficient group.  Finally, surjectivity of $q$
makes $S\to S\otimes_M H$ surjective: every class $[s,h]$
has the form $[s,q(m)]=[sm,1]$.  This proves the state bound.
\end{proof}

\begin{corollary}[Proved: transferring a finite-field square obstruction]
\label{CP2:A-band-quotient-modtwo-square-transfer}
In the preceding situation, if every degree-two class of
the target collage has square zero over $\mathbb F_2$, then
the source collage cannot have the homotopy type of
$\mathbb CP^2$.
For a hypothetical projective-plane source, the target has
$H_1(-;\mathbb F_2)=0$ and
$\dim H_2(-;\mathbb F_2)=1$.
\end{corollary}
\begin{proof}
Every degree-two class of the source is the pullback of a
target class, and the actual functor preserves cup products.
The nonzero square of the projective-plane generator over
$\mathbb F_2$ is therefore impossible.  The homology
assertions follow from the theorem and the natural
homological universal-coefficient sequence.
\end{proof}

This comparison puts no bound on the higher homology of
the target.  In a finite calculation used through this
corollary, a positive third Betti number of the target is
therefore not an exclusion by itself.  Likewise an actual
proper retract of the target can be omitted only with a
separate zero-square argument for that retract; the
projective-plane lower bound for a smaller category is
not, by itself, such an argument.
